\subsection{Matrix calculations.}

In this no., we propose to relax the matrix calculus introduced in Chap. II, § 6, and to apply it to translate certain results proved in this section.

I. --- Let I and K be two finite index sets, H a nonempty set, and \(M = (m_{ik})_{(i,k) \in \mathbb{I} \times \mathbb{K}}\) a matrix over H (chap. II, § 6, no. 1, def. 1).

The transpose of \(M\) is called, and denoted \({}^t M\) , the matrix \((m_{ki}')_{(k,i) \in \mathbb{K} \times \mathbf{I}}\) satisfying \(m_{ki}' = m_{ik} ((i, k) \in \mathbf{I} \times \mathbf{K})\) . We obviously have

\[
^ {t} (^ {t} M) = M.\tag{39}
\]

This generalizes the notion introduced in Chap. II, § 6, no. 6.

Suppose that H is a commutative group (written additively). The set of matrices over H with I and K as index sets admits a structure of commutative group, since it is the set of maps from I × K to H. This group is written additively.

Let H', H'' be two nonempty sets, H an abelian group (written additively), and \(f:(h',h'')\to h'h''\) a map of \(\mathbf{H}'\times\mathbf{H}''\) into H. Given two matrices

\[
M ^ {\prime} = (m _ {i k} ^ {\prime}) _ {(i, k) \in \mathrm{I} \times \mathrm{K}}, \quad M ^ {\prime \prime} = (m _ {k l} ^ {\prime \prime}) _ {(k, l) \in \mathrm{K} \times \mathrm{L}}
\]

on H' and H'' respectively, such that the set K of column indices of M' is equal to the set of row indices of M'' , one calls the product of M' and M'' (following f) and denotes it M'M'' the matrix

\[
M ^ {\prime}. M ^ {\prime \prime} = (\sum_ {k \in \mathbf {K}} m _ {i k} ^ {\prime} m _ {k l} ^ {\prime \prime}) _ {(i, l) \in \mathbf {I} \times \mathbf {L}}\tag{40}
\]

on H. This generalizes the notion introduced in chap. II, § 6, no. 4. If \(\mathbf{H}' = \mathbf{H}'' = \mathbf{H}\) and if H is a ring, the product \(M'M''\) will, unless expressly stated otherwise, be computed "in H", that is, according to the map \((x, y) \to xy\) . When \(\mathbf{H}'\) and \(\mathbf{H}''\) are commutative groups (written additively) and \(f\) is bilinear, one has

\[
\left\{ \begin{array}{l} (M ^ {\prime} + M _ {1} ^ {\prime}) M ^ {\prime \prime} = M ^ {\prime} M ^ {\prime \prime} + M _ {1} ^ {\prime} M ^ {\prime \prime}, \\ M ^ {\prime} (M ^ {\prime \prime} + M _ {1} ^ {\prime \prime}) = M ^ {\prime} M ^ {\prime \prime} + M ^ {\prime} M _ {1} ^ {\prime \prime}, \end{array} \right.\tag{41}
\]

where \(M'\) , \(M_1'\) are matrices over \(H'\) , \(M''\) , \(M_1''\) matrices over \(H''\) , and where the sums and products written are assumed to be defined. Let \(M'\) , \(M''\) be matrices over the sets \(H'\) , \(H''\) , and \(f^0\) the map of \(H'' \times H'\) into H defined by \((h'', h') \to h' h''\) ; then one has

\[
{ } ^ { t } ( M ^ { \prime } M ^ { \prime \prime } ) = { } ^ { t } M ^ { \prime \prime } . { } ^ { t } M ^ { \prime }\tag{42}
\]

where the product in the first (resp. second) member is computed according to \(f\) (resp. \(f^{0}\) ).

In the case where \(H' = H'' = H\) is a ring, one recovers formula (12) of chap. II, §6, no. 6.

Let A be a ring, J an antiautomorphism of A. For every matrix \(M = (m_{ik})\) over A, we shall denote by \(M^{J}\) the matrix \((m_{ik}^{J})\) . Let \(M_{1}\) , \(M_{2}\) be two matrices over A such that \(M_{1}M_{2}\) is defined. Since J is an isomorphism of A onto the opposite ring \(A^{0}\) , one has \((M_{1}M_{2})^{J} = M_{1}^{J}\) . \(M_{2}^{J}\) where the first (resp. second) member is computed in A (resp. \(A^{0}\) ). In view of (42) and (39), this gives

\[
(M _ {1} M _ {2}) ^ {\mathrm{J}} = ^ {t} (^ {t} M _ {2} ^ {\mathrm{J}}. ^ {t} M _ {1} ^ {\mathrm{J}})\tag{43}
\]

where both members are computed in A.

Let \(H_{1}\) , \(H_{2}\) , \(H_{3}\) , \(H_{12}\) , \(H_{23}\) and H be commutative groups (written additively), \(f_{12}: H_{1} \times H_{2} \to H_{12}\) , \(f_{23}: H_{2} \times H_{3} \to H_{23}\) , \(f_{3}: H_{12} \times H_{3} \to H\) , \(f_{1}: H_{1} \times H_{23} \to H\) maps, and let \(M_{1}\) , \(M_{2}\) , \(M_{3}\) be matrices over \(H_{1}\) , \(H_{2}\) , \(H_{3}\) respectively. If \(f_{3}(f_{12}(x_{1}, x_{2}), x_{3}) = f_{1}(x_{1}, f_{23}(x_{2}, x_{3}))\) for all \(x_{i} \in H_{i}\) (i = 1, 2, 3), then the products \((M_{1}M_{2})M_{3}\) and \(M_{1}(M_{2}M_{3})\) (computed according to \(f_{12}\) , \(f_{3}\) , \(f_{23}\) and \(f_{1}\) ), if they are defined, are equal; they will be denoted \(M_{1}M_{2}M_{3}\) . When \(H_{1} = H_{2} = H_{3} = H_{12} = H_{23} = H\) , that H is a ring, and that \(f_{12}\) , \(f_{23}\) , \(f_{3}\) , \(f_{1}\) are equal to the map \((x, y) \to xy\) , the preceding condition expresses the associativity of the latter, and is therefore satisfied. Analogous conventions will be made for products of more than three factors.

Let A, B be two rings, \(M = (m_{ik})_{(i,k) \in I \times K}\) and \(M' = (m_{ik}')_{(i,k) \in I \times K}\) be two matrices over an (A, B)-bimodule G (no. 1). If, for every one-row matrix \(L = (a_i)_{i \in I}\) with elements in A and every one-column matrix \(C = (b_k)_{k \in K}\) with elements in B, one has \(L.M.C = L.M'.C\) (the products being computed according to the maps defining the structure of (A, B)-bimodule of G), then the matrices \(M\) and \(M'\) are equal. Indeed, if one takes \(a_i = 1\) , \(a_s = 0\) for \(s \neq i\) , \(b_k = 1\) , \(b_t = 0\) for \(t \neq k\) , the matrices \(L.M.C\) and \(L.M'.C\) , which are scalar matrices, are respectively equal to \(m_{ik}\) and \(m_{ik}'\) .

\begin{enumerate}
\def\labelenumi{\Roman{enumi}.}
\setcounter{enumi}{1}
\tightlist
\item
  Consider a ring A and a (right or left) A-module E, admitting a finite basis \((e_{i})_{i\in I}\) . For every element x of E, the one-column matrix formed by the components \(x_{i}\) ( \(i\in I\) ) of x with respect to \((e_{i})\) is called the matrix of x with respect to the basis \((e_{i})\) , and is denoted \(M(x)\) or x (cf. chap. II, § 6, no. 4); in the calculations it will be convenient, in order to recall that the index i is a row index, to adjoin to it a column index capable of taking only one value, and to write \((x_{i0})\) for the matrix \(M(x)\) .
\end{enumerate}

Now consider two A-modules (left or right) E and F, having finite bases \((e_{i})_{i\in I}\) and \((f_{k})_{k\in K}\) respectively; let \((f_{k}^{*})\) be the basis of \(F^{*}\) dual to \((f_{k})\) . We shall define the matrix, with respect to these bases, of a map u from E to F in the following four cases:

(D) E and F are right modules, and u is A-linear;

(G) E and F are left modules, and u is A-linear;

(GD) E is a left module, F a right module, A is equipped with an antiautomorphism J, u is Z-linear and satisfies \(u(ax) = u(x)a^{J}\) ( \(a \in A, x \in E\) ) (in other words, u is an A-linear map from \(E^{J}\) to F (no. 2, def. 5)).

(DG) E is a right A-module, F a left A-module, A is equipped with an antiautomorphism J, u is Z-linear and satisfies \(u(xa) = a^{J}u(x)\) ( \(x \in E, a \in A\) ) (in other words, u is an A-linear map from \(E^{J}\) into F).

In each of these four cases, the matrix of the map u is, by definition, the matrix \((u_{ki})_{(k,i)\in\mathbb{K}\times\mathbb{I}}\) such that

\[
u _ {k i} = \left\langle u (e _ {i}), f _ {k} ^ {*} \right\rangle .\tag{44}
\]

This definition coincides, in case (D), with that given in Chap. II, § 6, no. 3. Under these conditions, the matrix \(M(u(x))\) of the image of an element \(x\) of E is given by the following formulas:

\[
M (u (x)) = M (u). M (x)\tag{45 D}
\]

\[
{ } ^ { t } M ( u ( x ) ) = { } ^ { t } M ( x ) . { } ^ { t } M ( u )\tag{45 G}
\]

\[
M (u (x)) = M (u). M (x) ^ {\mathrm{J}}\tag{45 GD}
\]

\[
{ } ^ { t } M ( u ( x ) ) = { } ^ { t } M ( x ) ^ { j } . { } ^ { t } M ( u ) .\tag{45 DG}
\]

Let us verify, for example, (45 DG), the other verifications being analogous and slightly easier. Set \(x = \sum_{i} e_i x_{io}, u(x) = \sum_{k} y_{ko} f_k\) ; we have \(u(x) = u(\sum_{i} e_i x_{io}) = \sum_{i} x_{io}^{\mathbf{J}} u(e_i) = \sum_{i,k} x_{io}^{\mathbf{J}} u_{ki} f_k\) ; whence \(y_{ko} = \sum_{i} x_{io}^{\mathbf{J}} u_{ki}\) ; in order to place the two indices \(i\) next to each other, consider the transposed matrices \(^t M(x) = (x_{oi}')\) where \(x_{oi}' = x_{io}\) , and \(^t M(u) = (u_{ik}')\) where \(u_{ik}' = u_{ki}\) ; we then have \(y_{ko} = \sum_{i} x_{oi}' u_{ik}'\) ; since the right-hand side is the element with index \(k\) of the one-row matrix \(^t M(x)^{\mathbf{J}}\) . \(^t M(u)\) , formula (45 DG) is verified.

Remarks. -- 1) When A is commutative, (45 G) reduces to (45 D), and (45 DG) to (45 GD), by means of the formula \(t(M'M'') = 'M'''.'M'\) (cf. (42)), where both sides are here computed in A. 2) Let E, F, G be three left modules having finite bases, and \(u: \mathrm{E} \to \mathrm{F}\) , \(\nu: \mathrm{F} \to \mathrm{G}\) A-linear maps. It follows from (45 G) that one has

\[
{ } ^ { t } M ( \nu \circ u ) = { } ^ { t } M ( u ) . { } ^ { t } M ( \nu ) .\tag{46}
\]

Indeed, we have, for any \(x \in E\) ,

\[
\begin{array}{r l} ^ {t} M (x). ^ {t} M (\nu \circ u) & = ^ {t} M (\nu (u (x))) = ^ {t} M (u (x)). ^ {t} M (\nu) \\ & = ^ {t} M (x). ^ {t} M (u). ^ {t} M (\nu), \end{array}
\]

whence (46).

Recall that, in the case of right modules, we have

\[
M (\nu \circ u) = M (\nu) M (u).
\]

\begin{enumerate}
\def\labelenumi{\Roman{enumi}.}
\setcounter{enumi}{2}
\tightlist
\item
  From now on, we denote by A a ring, by B a ring (resp. a ring equipped with an antiautomorphism J, for which we set \(J' = J^{-1}\) ), by E a left A-module having a finite basis \((e_i)_{i \in I}\) and by F a right (resp. left) B-module having a finite basis \((f_k)_{k \in K}\) . We denote by \((e_i^*)\) and \((f_k^*)\) the dual bases of \(E^*\) and \(F^*\) . Unless expressly stated otherwise, the matrices considered are taken with respect to these bases.
\end{enumerate}

Let G be an (A, B)-bimodule (no. 1), Φ a bilinear (resp. sesquilinear on the right for J) map of E × F into G, and \(R = (\Phi(e_i, f_k))\) the matrix of Φ. Then, for \(x \in E\) and \(y \in F\) , formula (6) of no. 1 (resp. (8) of no. 2) is written, by virtue of the above conventions,

\[
\Phi (x, y) = ^ {t} M (x). R. M (y) \quad (\text { resp. } \Phi (x, y) = ^ {t} M (x). R. M (y) ^ {\mathrm{J}}),\tag{47}
\]

where the products are computed according to the maps which define the structure of the (A, B)-bimodule G; in particular, if A = B = G (in which case Φ is a form), the products are computed in A.

Let E' be a left A-module having a finite basis \((e_{s}^{\prime})_{s\in S}\) , F' a right (resp. left) A-module having a finite basis \((f_{t}^{\prime})_{t\in T}\) , \(u:E\to E^{\prime}\) and \(\nu:F\to F^{\prime}\) A-linear maps, and \(\Phi^{\prime}\) a bilinear map (resp. sesquilinear on the right for J) from \(E^{\prime}\times F^{\prime}\) into G. Denote by \(\Phi\) the inverse image of \(\Phi^{\prime}\) (relative to u and \(\nu\) ), U, V, R, \(R^{\prime}\) the matrices of u, \(\nu\) , \(\Phi\) , \(\Phi^{\prime}\) with respect to the bases considered. One then has

\[
R = ^ {t} U. R ^ {\prime}. V \quad (\text { resp. } R = ^ {t} U. R ^ {\prime}. V ^ {J}),\tag{48}
\]

the products being computed as in (47). Indeed, whatever \(x \in \mathbf{E}\) and \(y \in \mathbf{F}\) , we have by definition \(\Phi(x, y) = \Phi'(u(x), v(y))\) , whence, by (47),

\[
\begin{array}{c} ^ {t} M (x). R. M (y) = ^ {t} M (u (x)). R ^ {\prime}. M (\nu (y)) \\ (\text {resp.} ^ {t} M (x). R. M (y) ^ {\mathrm{J}} = ^ {t} M (u (x)). R ^ {\prime}. M (\nu (y)) ^ {\mathrm{J}}); \end{array}
\]

by (45 G) and (45 D) (resp. (45 G)) and (43) one deduces

\[
\begin{array}{r l} ^ {t} M (x). R. M (y) & = ^ {t} M (x). ^ {t} U. R ^ {\prime}. V. M (y) \\ (\text {resp.} ^ {t} M (x). R. M (y) ^ {J} & = ^ {t} M (x). ^ {t} U. R ^ {\prime}. ^ {t} (^ {t} M (y). ^ {t} V) ^ {J} \\ & = ^ {t} M (x). ^ {t} U. R ^ {\prime}. V ^ {J}. M (y) ^ {J}); \end{array}
\]

this proves our assertion.

\begin{enumerate}
\def\labelenumi{\Roman{enumi}.}
\setcounter{enumi}{3}
\tightlist
\item
  --- We assume here that the rings A and B are equal, and we denote by \(\Phi\) a bilinear (resp. sesquilinear on the right for J) form on E × F, and by \(R\) its matrix. Let us compute the matrices of the maps \(s_{\Phi}\) and \(d_{\Phi}\) associated with \(\Phi\) , which we shall denote by \(s\) and \(d\) to simplify. Since one has \(\Phi(x, y) = \langle y, s(x) \rangle = \langle x, d(y) \rangle\) by (23), no. 6 (resp. \(\Phi(x, y) = \langle x, d(y) \rangle = \langle y, s(x) \rangle^{J}\) by (24), no. 6), one has \(\Phi(e_i, f_k) = \langle f_k, s(e_i) \rangle = \langle e_i, d(f_k) \rangle\) (resp. \(\Phi(e_i, f_k) = \langle e_i, d(f_k) \rangle = \langle f_k, s(e_i) \rangle^{J}\) ), whence, by (44) and since ( \(e_i\) ) is the dual basis of ( \(e_i^*\) ) and ( \(f_k\) ) the dual basis of ( \(f_k^*\) ):
\end{enumerate}

\[
M (d) = R,   M (s) = ^ {t} R \quad (\text { resp. }   M (d) = R,   M (s) = ^ {t} R ^ {J ^ {\prime}}).\tag{49}
\]

Remarks. --- 1) When A is a field, the linear maps s and d have the same rank. We see here that their matrices \(M(s)\) and \(M(d)\) have the same rank; indeed, a matrix over A and its transpose have the same rank (chap. II, § 6, no. 7, prop. 3) and, when \(\Phi\) is sesquilinear, the equality of the ranks of \(R\) over A and of \(^tR\) over A \(^0\) (ibid.) and the fact that J' is an isomorphism of A \(^0\) onto A, imply the equality of the ranks of \(R\) and of \(^tR^{J'}\) over A.

\begin{enumerate}
\def\labelenumi{\arabic{enumi})}
\setcounter{enumi}{1}
\tightlist
\item
  If M and N are right A-modules having finite bases \((m_{i})\) and \((n_{k})\) , \(\Phi\) a sesquilinear form on the left for J on \(M \times N\) (no. 6, Remark), s and d its associated maps, and \(R = (\Phi(m_{i}, n_{k}))\) its matrix, formulas (26) of no. 6 show that one has
\end{enumerate}

\[
M (d) = R ^ {J ^ {\prime}}, \quad M (s) = ^ {t} R.
\]

Suppose now that the maps s and d associated with \(\Phi\) are bijective and let us compute the matrix \(\hat{R}\) of the inverse form of \(\Phi\) (no. 7). When \(\Phi\) is bilinear, \(\Phi\) is the inverse image of \(\hat{\Phi}\) with respect to the linear maps \(s: \mathrm{E} \to \mathrm{F}^*\) and \(d: \mathrm{F} \to \mathrm{E}^*\) ; one therefore has, by virtue of (48) and (49), \(R = R\) . \(\hat{R}\) . \(R\) , whence, since \(R\) is invertible ( \(d\) being bijective), \(\hat{R} = R^{-1}\) . This formula extends to the case where \(\Phi\) is sesquilinear, for, if one considers \(\Phi\) as a bilinear form on \(\mathrm{E} \times \mathrm{F}^{\mathrm{j}}\) , and if one identifies \((\mathrm{F}^{\mathrm{j}})^{*}\) with \((\mathrm{F}^{*})^{\mathrm{j}}\) (cf. no. 9), the inverse form of this bilinear form coincides with \(\hat{\Phi}\) considered as a bilinear form on \((\mathrm{F}^{*})^{\mathrm{j}} \times \mathrm{E}^{*}\) . In both cases the matrix of the inverse form of \(\Phi\) is the inverse of the matrix of \(\Phi\) .

Finally, let E' be a left A-module, F' a right (resp. left) A-module, both admitting finite bases \((e_{s}^{\prime})\) and \((f_{t}^{\prime})\) ; let \(\Phi^{\prime}\) be a bilinear form (resp. sesquilinear form with respect to J) on \(E^{\prime} \times F^{\prime}\) , and let \(R^{\prime}\) be its matrix. Suppose \(s_{\Phi}\) and \(d_{\Phi}\) bijective. Let \(u : E \to E^{\prime}\) and \(\nu : F \to F^{\prime}\) be linear maps, \(u^{*} : F^{\prime} \to F\) and \(\upsilon^{*} : E^{\prime} \to E\) their adjoints (no. 8, prop. 7); denote by U, V, \(U^{*}\) , \(V^{*}\) the matrices of u, \(\upsilon\) , \(u^{*}\) , \(\upsilon^{*}\) with respect to the given bases. One then has

\[
U ^ {*} = R ^ {- 1}. ^ {t} U. R ^ {\prime}, \quad {} ^ {t} V ^ {*} = R ^ {\prime}. V. R ^ {- 1}\tag{50}
\]

\[
(\text { resp. } U ^ {* \mathrm{J}} = R ^ {- 1}. ^ {t} U. R ^ {\prime}, ^ {t} V ^ {*} = R ^ {\prime}. V ^ {\mathrm{J}}. R ^ {- 1}).
\]

Indeed, whatever \(x \in E\) and \(y \in F'\) , one has \(\Phi'(u(x), y) = \Phi(x, u^*(y))\) (no. 8, def. 10). Hence, when \(\Phi\) is bilinear, by virtue of (47), \(^t M(u(x)). R'. M(y) = ^t M(x). R. M(u^*(y))\) ; this gives, by virtue of (45 G) and (45 D), \(^t M(x). ^t U. R'. M(y) = ^t M(x). R. U^*. M(y)\) , whence \(^t U. R' = R. U^*\) and the first announced formula, since, \(d\) being bijective, \(R\) is invertible. When \(\Phi\) is sesquilinear (47) and (45 G) give \(^t M(x). ^t U. R'. M(y)^J = ^t M(x). R. (^{t}M(y). ^{t}U^*)^{J}\) ; now, according to (43), we have \((^{t}M(y). ^{t}U^*)^{J} = ^{t}(^{u}U^{*J}. ^{u}M(y)^{J})\) , whence \((^{t}M(y). ^{t}U^*)^{J} = U^{*J}. M(y)^{J}\) ; it follows therefore \(^t M(x). ^{t}U. R'. M(y)^{J} = ^{t}M(x). R. U^{*J}. M(y)^{J}\) , whence \(^t U. R' = R. U^{*J}\) , and \(U^{*J} = R^{-1}. ^{t}U. R'\) . The verification of the formulas for \(V^*\) is analogous.

Exercises. --- 1) Let A be a commutative field, E a vector space over A admitting an infinite countable basis \((e_n)_{n \geqslant 1}\) . We define a bilinear form \(\Phi\) on E by setting \(\Phi(e_{i+1}, e_i) = 1\) for \(i \geqslant 1\) , \(\Phi(e_k, e_j) = 0\) for \(k \neq j + 1\) and \(j \geqslant 1\) . Show that the linear map \(d_\Phi\) associated on the right with \(\Phi\) is injective, but the linear map \(s_\Phi\) associated on the left with \(\Phi\) is not injective.

\begin{enumerate}
\def\labelenumi{\arabic{enumi})}
\setcounter{enumi}{1}
\item
  Let E be the Z-module direct sum of Z and Z/(2), and let E* be its dual (isomorphic to Z). Show that the bilinear form \((x, x') \to \langle x, x' \rangle\) on \(E \times E^*\) is such that the associated linear map on the right is injective, but the associated linear map on the left is not.
\item
  Give an example of a bilinear form \(\Phi\) defined on a product \(\mathbf{E} \times \mathbf{F}\) of two vector spaces, such that \(d_{\Phi}\) is bijective, \(s_{\Phi}\) injective but not bijective (take \(\mathbf{E}\) of infinite dimension and F equal to the dual \(\mathbf{E}^*\) of \(\mathbf{E}\) ; cf. chap. II, § 5, exerc. 3).
\item
  Let A be a ring equipped with an antiautomorphism J, let E be a left A-module, and let G be an (A, A)-bimodule, and \(\Phi\) a map from \(\mathrm{E} \times \mathrm{E}\) into G, sesquilinear on the right for J. Prove the identity (where \(\mathrm{Q}(x) = \Phi(x, x)\) ):
\end{enumerate}

\[
\begin{array}{r l} & {2 \Phi (x, y) (\mu^ {\mathrm{J}} \lambda^ {\mathrm{J}} - \lambda^ {\mathrm{J}} \mu^ {\mathrm{J}}) = \mathrm{Q} (x - \mu \lambda y) - \mathrm{Q} (x + \mu \lambda y) + \mu \mathrm{Q} (x + \lambda y)} \\ & {- \mu \mathrm{Q} (x - \lambda y) + \mathrm{Q} (x + \mu y) \lambda^ {\mathrm{J}} - \mathrm{Q} (x - \mu y) \lambda^ {\mathrm{J}} + \mu \mathrm{Q} (x - y) \lambda^ {\mathrm{J}} - \mu \mathrm{Q} (x + y) \lambda^ {\mathrm{J}}.} \end{array}
\]

\begin{enumerate}
\def\labelenumi{\arabic{enumi})}
\setcounter{enumi}{4}
\tightlist
\item
  Let K be a commutative field of characteristic 2, A a separable quadratic extension of K; we have \(\mathrm{A} = \mathrm{K}(\theta)\) , where \(\theta\) is a root of an irreducible polynomial \(\mathrm{X}^2 +\mathrm{X} + \beta\) of K{[}X{]} and the K-automorphism J of A, distinct from the identity, is such that \(\theta^{j} = \theta +1\) (chap. V, § 11, exerc. 8). Show that if E and G are vector spaces over A, \(\Phi\) a sesquilinear map (for J) from \(\mathrm{E}\times \mathrm{E}\) into G, we have, setting \(\mathrm{Q}(x) = \Phi (x,x)\) ,
\end{enumerate}

\[
\Phi (x, y) = \mathrm{Q} (\theta x + y) - \beta \mathrm{Q} (x) - \mathrm{Q} (y) - (\theta + 1) (\mathrm{Q} (x + y) - \mathrm{Q} (x) - \mathrm{Q} (y)).
\]

\begin{enumerate}
\def\labelenumi{\arabic{enumi})}
\setcounter{enumi}{5}
\tightlist
\item
  Let A be a field, E a vector space over A, \(\Phi\) a sesquilinear form on E, \(u\) an endomorphism of E.
\end{enumerate}

\begin{enumerate}
\def\labelenumi{\alph{enumi})}
\item
  For there to exist one and only one endomorphism \(u^*\) of E such that \(\Phi(u(x), y) = \Phi(x, u^*(y))\) for \(x, y\) in E, it is necessary and sufficient that \(d_\Phi\) be injective and that \(^t u(d_\Phi(E)) \subset d_\Phi(E)\) .
\item
  Give an example where E is infinite-dimensional and \(d_{\Phi}\) injective, but where \(u(d_{\Phi}(E))\) is not contained in \(d_{\Phi}(E)\) .
\end{enumerate}

\begin{enumerate}
\def\labelenumi{\arabic{enumi})}
\setcounter{enumi}{6}
\tightlist
\item
  Let E, \(E_{1}\) be two A-modules, \(\Phi\) (resp. \(\Phi_{1}\) ) a sesquilinear form on E (resp. \(E_{1}\) ). Assume that \(\Phi_{1}\) is nondegenerate and that there exists an element \(\alpha\in A\) and a bijection u from E onto \(E_{1}\) such that \(\Phi_{1}(u(x), u(y))=\Phi(x,y)\alpha\) for all x, y in E. Show that: \(1^{\circ}\Phi\) is nondegenerate; \(2^{\circ}u\) is linear; \(3^{\circ}\) if \(E_{1}\) is a faithful A-module, so is E, and \(\alpha\) is not a right zero divisor in A; \(4^{\circ}\) if \(\Phi_{1}\) takes values in A that are not left zero divisors, so is \(\Phi\) .
\end{enumerate}

¶ 8) Let A be a field, E₁, E₂ two nonzero vector spaces over A, Φ₁ (resp. Φ₂) a nondegenerate sesquilinear form on E₁ (resp. E₂) for an antiautomorphism J₁ (resp. J₂) of A. Let u be a linear map from E₁ onto E₂ such that the relation Φ₁(x, y) = 0 implies Φ₂(u(x), u(y)) = 0.

\begin{enumerate}
\def\labelenumi{\alph{enumi})}
\item
  Show that \(u\) is a bijection from \(\mathbf{E}_{1}\) onto \(\mathbf{E}_{2}\) . (If \(u(0)\) were not reduced to 0, show that there would exist in \(\mathbf{E}_{1}\) two vectors \(a\) , \(b\) such that \(u(a) \neq 0\) , \(u(b) = 0\) and \(\Phi_{1}(a, b) \neq 0\) ; if H is the hyperplane of \(x \in \mathbf{E}_{1}\) such that \(\Phi_{1}(a, x) = 0\) , note that one would have \(u(\mathrm{H}) = \mathrm{E}_{2}\) ).
\item
  Show that if dim \(E_{1} \geqslant 2\) , there exists \(\alpha \in A\) such that one has \(\Phi_{2}(u(x), u(y)) = \Phi_{1}(x, y) \alpha\) for all x, y in \(E_{1}\) . (For every \(y \in E_{1}\) , show that there exists an element \(m(y) \in A\) such that \(\Phi_{2}(u(x), u(y)) = \Phi_{1}(x, y)m(y)\) for all \(x \in E_{1}\) , and that if y and \(y'\) are linearly independent in \(E_{1}\) , one has \(m(y + y') = m(y) = m(y')\) ).
\end{enumerate}

\begin{enumerate}
\def\labelenumi{\arabic{enumi})}
\setcounter{enumi}{8}
\tightlist
\item
  Let A be a field, E, F be two left vector spaces over A, \(\Phi\) a nondegenerate sesquilinear form on \(E \times F\) for an antiautomorphism J of A.
\end{enumerate}

\begin{enumerate}
\def\labelenumi{\alph{enumi})}
\item
  Let M be a subspace of E, and N a subspace of F such that \(N \supset M^{0}\) and \(M \supset N^{0}\) . Show that if one of the spaces \(N/M^{0}\) , \(M/N^{0}\) is finite-dimensional, then so is the other, and the dimensions of these two spaces are equal.
\item
  Let M, \(M'\) be two subspaces of E such that \(M^{00} = M\) and \(M'\) be of finite dimension; show that we have \((M \cap M')^{0} = M^{0} + M'^{0}\) and \((M + M')^{00} = M + M'\) . (Applying a) to the subspaces \(M'\) and \(M^{0} + M'^{0}\) , show that \(\dim(M \cap M') = \text{codim}(M^{0} + M'^{0})\) ; applying a) to the subspaces \(M + M'\) and \(M^{0}\) , show that
\end{enumerate}

\[
\dim ((M + M ^ {\prime}) ^ {0 0} / M) = \dim ((M + M ^ {\prime}) / M).
\]

\begin{enumerate}
\def\labelenumi{\alph{enumi})}
\setcounter{enumi}{2}
\item
  If E = F and if M is a subspace of E such that \(E = M^{0} + M^{00}\) , show that E is the direct sum of \(M^{0}\) and \(M^{00}\) .
\item
  Let E be a vector space over a commutative field A admitting an infinite countable basis \((e_{n})_{n\geqslant0}\) , and let \(\Phi\) be the symmetric bilinear form on E such that \(\Phi(e_{n}, e_{n}) = 1\) for all n, \(\Phi(e_{i}, e_{j}) = 0\) for \(i \geqslant 1\) , \(j \geqslant 1\) and \(i \neq j\) , and \(\Phi(e_{0}, e_{n}) = 1\) for all \(n \geqslant 1\) . Show that \(\Phi\) is nondegenerate. Let M (resp. N) be the subspace of E spanned by the \(e_{2k}\) (resp. \(e_{2k-1}\) ) for \(k \geqslant 1\) , and let H = M + N, which is a hyperplane in E. Show that we have \(M^{0} = N\) , \(N^{0} = M\) , \(H^{00} = E \neq H\) , \((M \cap N)^{0} \neq M^{0} + N^{0}\) and \((M + N)^{00} \neq M + N\) , although \(M^{00} = M\) , \(N^{00} = N\) ; if L is the 2-dimensional subspace spanned by \(e_{0}\) and \(e_{1}\) , one has \((L \cap H)^{0} \neq L^{0} + H^{0}\) .
\end{enumerate}

¶ 10) Let E, E' be two left vector spaces over fields A, A' respectively, of dimension ≥3; let ℱ(E) (resp. ℱ(E')) be the lattice-ordered set (for the inclusion relation) formed by the finite-dimensional subspaces of E (resp. E').

\begin{enumerate}
\def\labelenumi{\alph{enumi})}
\item
  Let \(p\) be a map from \(\mathfrak{F}(\mathrm{E})\) into \(\mathfrak{F}(\mathrm{E}')\) such that for every \(\mathbf{M} \in \mathfrak{F}(\mathrm{E})\) , \(\dim p(\mathbf{M}) = \dim \mathbf{M}\) , and such that for every pair \((\mathbf{M}, \mathbf{N})\) of elements of \(\mathfrak{F}(\mathrm{E})\) , \(p(\mathbf{M} + \mathbf{N}) = p(\mathbf{M}) + p(\mathbf{N})\) . Show that \(p\) is injective; if \(p\) is bijective, there exists a bijective semilinear map \(u\) of \(\mathrm{E}\) into \(\mathrm{E}'\) such that one has \(u(\mathbf{M}) = p(\mathbf{M})\) for all \(\mathbf{M} \in \mathfrak{F}(\mathrm{E})\) (use exercise 10 of chap. II, 2nd ed., App. III).
\item
  Give an example where \(A' = A\) is commutative, \(E' = E\) is of finite dimension, and where there exists a map p from \(\mathfrak{F}(E)\) into itself, such that \(\dim p(M) = \dim M\) , \(p(M + N) = p(M) + p(N)\) , \(p(M \cap N) = p(M) \cap p(N)\) , but there exists no injective semilinear map u from E into itself such that \(u(M) = p(M)\) for \(M \in \mathfrak{F}(E)\) . (Consider the case where there exists an overfield \(A''\) of A of finite degree and isomorphic to A, for example \(\mathbf{A} = \mathbf{F}_{p}(\mathbf{X})\) , where \(p\) is prime; \(\mathrm{E}'' = \mathrm{A}'' \otimes_{\mathbf{A}} \mathrm{E}\) is then a vector space of the same dimension over \(\mathrm{A}''\) as \(\mathrm{E}\) over \(\mathrm{A}\) ; consider the map \(\mathrm{M} \to \mathrm{A}'' \otimes_{\mathbf{A}} \mathrm{M}\) .
\item
  We assume \(A' = A\) . Let \(\omega\) be a bijective map from \(\mathfrak{F}(E)\) onto the set \(\mathfrak{F}'(E')\) of subspaces of finite codimension of \(E'\) , such that \(\text{codim } \omega(M) = \dim M\) , and \(\omega(M + N) = \omega(M) \cap \omega(N)\) . Show that there exists a sesquilinear form \(\Phi\) on \(E \times E'\) , nondegenerate on the right, and such that \(\omega(M) = M^0\) for all \(M \in \mathfrak{F}(E)\) . (Use th. 1 of chap. II, § 4, no. 6).
\end{enumerate}

¶ 11) a) Let A be a ring artinian on the left and on the right (chap. VIII, § 2, no. 3). Show that the following conditions are equivalent: 1° the right annihilator (resp. left annihilator) of every left ideal (resp. right ideal) ≠ A is not reduced to 0; 2° the dual of every simple left A-module (resp. right A-module) is not reduced to 0; 3° the dual of every finitely generated left A-module (resp. right A-module) is not reduced to 0. A ring A is said to satisfy condition (N \(_{s}\) ) (resp. (N \(_{a}\) )) if it satisfies these conditions for left A-modules (resp. right A-modules).

\begin{enumerate}
\def\labelenumi{\alph{enumi})}
\setcounter{enumi}{1}
\item
  Let A be a ring artinian on the left and on the right satisfying condition (N \(_{a}\) ), E (resp. F) a free left A-module (resp. right A-module) having a countable basis over A, Φ a bilinear form on E × F such that s \(_{\Phi}\) is injective. Let M be a free submodule of E; show that there exists a free submodule N of F and a basis (e \(_{n}\) ) (resp. (f \(_{n}\) )) of M (resp. N) such that Φ(e \(_{i}\) , f \(_{j}\) ) = δ \(_{ij}\) ; moreover one can take for e \(_{1}\) any free element of M. (Note that if x is a free element of M, the image of F under the map y → Φ(x, y) is the whole ring A; then proceed by induction to construct the bases (e \(_{n}\) ) and (f \(_{n}\) )). Deduce that if the basis (e \(_{n}\) ) of M is finite, E is the direct sum of M and N \(^{0}\) , and that one has M \(^{00}\) = M.
\item
  We keep the hypotheses of \(b\) , and suppose in addition that A satisfies condition \((\mathrm{N}_{s})\) and that \(d_{\Phi}\) is injective. Show then that there exists a basis \((e_{n})\) in E and a basis \((f_{n})\) in F such that \(\Phi(e_{i}, f_{j}) = \delta_{ij}\) . (Use \(b\) ) by determining alternately by induction \(e_{n}\) and \(f_{n}\) ).
\end{enumerate}

¶ 12) Let E be a real Hilbert space of countable type, \(\Phi(x, y)\) the scalar product in E. Show that there does not exist a system of two algebraic bases \((e_{\lambda})\) , \((f_{\mu})\) of the vector space E over R such that we have \(\Phi(e_{\lambda}, f_{\mu}) = \delta_{\lambda y}\) for every pair of indices. (First note that the index set of these bases would have the cardinality of the continuum (Topological vector spaces, chap. II, § 3, exerc. 15); then consider an orthonormal (countable) basis \((a_n)\) of E and note that the subspace spanned by the \(a_n\) is contained in the subspace generated by a countable subfamily of \((e_{\lambda})\) .
% semantic-fragments: legacy-input-seam
\relax{}
